\section{Completion of the induction}\label{sec:completion}

We now combine the preceding constructions. The main geometric step is
to make the ordinary adjoint semiample after adding a sufficiently
positive line from the algebraic base, while preserving the map to that
base. The adjoint formula of Proposition~\ref{adj:base} then puts the
problem on a projective variety. We first isolate the negative-part
calculation needed to pull the projective answer back.

\subsection{Divisors above subsets of codimension at least two}

\begin{lemma}\label{finish:exceptional}
Let \(r:V\to S\) be a surjective morphism from a smooth compact K\"ahler
manifold to a normal projective variety. Let \(H\) be a nef rational
Cartier divisor on \(S\). If \(E\geq0\) is a real divisor on \(V\) whose
every component has image of codimension at least two in \(S\), then
\[
 N(r^*H+E)=E.
\]
\end{lemma}
\begin{proof}
Put \(n=\dim V\), \(\beta=c_1(r^*H)\), and
\(\alpha=\beta+\{E\}\). Nefness gives \(N(\alpha)\leq E\).
Suppose
\[
 U=E-N(\alpha)>0.
\]
Then \(\beta+\{U\}=Z(\alpha)\) is modified nef. Write
\(U=\sum_i u_iU_i\), with \(u_i>0\), and choose
\[
 s=\max_i\dim r(U_i)\leq\dim S-2.
\]
In particular \(0\leq s\leq n-2\). Let \(\eta\) be the pullback of an
ample class on \(S\), and let \(\omega\) be a K\"ahler class on \(V\).
Consider the symmetric bilinear form
\[
 Q(\gamma,\delta)=\int_V\gamma\,\delta\,\eta^s\omega^{n-s-2}
 \quad\text{on }H^{1,1}(V,\R).
\]
Since the image of each \(U_i\) has dimension at most \(s\), any pairing
on \(U_i\) with \(s+1\) classes from \(S\) vanishes. Therefore
\begin{equation}\label{finish:orthogonal}
 Q(\{U\},\eta)=0,\qquad Q(\{U\},\beta)=0.
\end{equation}
By Lemma~\ref{pre:negative}(3), \(Z(\alpha)\) restricts
pseudo-effectively to a resolution of \(U_i\). Pairing there with the
nef classes \(\eta^s\omega^{n-s-2}\) and summing over \(i\) gives
\[
 0\leq Q(\{U\},Z(\alpha))=Q(\{U\},\{U\}).
\]
On the other hand,
\[
 Q(\eta,\eta)=\int_V\eta^{s+2}\omega^{n-s-2}>0,
\]
because \(s+2\leq\dim S\). The mixed Hodge--Riemann relations
\cite[Theorem~A]{DinhNguyen}, applied with
\(\eta+\varepsilon\omega\) and then passed to the limit, show that \(Q\)
has at most one positive eigenvalue. Thus its restriction to
\(\eta^\perp\) is negative semidefinite. Equations
\eqref{finish:orthogonal} force \(Q(\{U\},\{U\})=0\).
A zero-square vector for a negative semidefinite form belongs to its
radical; decomposing every class into a multiple of \(\eta\) and an
element of \(\eta^\perp\) shows that \(\{U\}\) belongs to the radical of
the whole form \(Q\).

This contradicts
\[
 Q(\{U\},\omega)=
 \sum_i u_i\int_{U_i}\eta^s\omega^{n-s-1}>0.
\]
Indeed every term is nonnegative, and a component with image dimension
\(s\) gives a strictly positive integral on its smooth generic locus.
Consequently \(U=0\).
\end{proof}

This proof is the codimension-two part of the lifting argument in
\cite[Section~5]{K}. Its formulation above also applies to maps with
positive-dimensional fibers; birational exceptional translation alone
would not suffice at that point of our proof.

\subsection{Keeping the ordinary program over the algebraic base}

\begin{proposition}\label{finish:over-base}
Let \(T\) be a smooth compact K\"ahler manifold of dimension \(n\), let
\((T,B)\) be a rational simple normal crossing klt pair with
\(J=K_T+B\) pseudo-effective, and let
\(h:T\to W\) be a surjective morphism to a smooth projective variety.
There are a smooth compact K\"ahler modification \(T_1\to T\), an
effective klt adjoint \(J_1\) on \(T_1\) as in
\eqref{pre:resolution-error}, and a finite ordinary \(J_1\)-negative
program over \(W\), ending at a normal compact K\"ahler klt pair
\((Y,\Delta)\), such that for some positive integer \(b\) and an ample
Cartier divisor \(A_W\) on \(W\),
\[
 K_Y+\Delta+b h_Y^*A_W
\]
is semiample. Here \(h_Y:Y\to W\) is the descended morphism.
\end{proposition}
\begin{proof}
Choose \(A_W\) very ample and an integer \(b>2n\), and set
\(C=h^*A_W\). A general member of a sufficiently high multiple of
the free system \(|C|\), divided by that multiple and multiplied by
\(b\), gives an effective rational representative of \(bC\).
Bertini's theorem, with the multiple chosen large enough, makes its
sum with \(B\) an effective simple normal crossing klt boundary.
Theorem~\ref{pre:ordinary} therefore gives a decomposition of \(J+bC\).

Pass to a higher smooth model \(T_1\) and apply the convention
\eqref{pre:resolution-error} to the original pair \((T,B)\). Write
\(J_1=K_{T_1}+B_1\) and \(C_1=h_1^*A_W\). The exceptional error and
Lemma~\ref{pre:negative} retain an actual identity
\begin{equation}\label{finish:augmented-decomposition}
 J_1+bC_1\qsim P_1+R_1,\qquad
 P_1\text{ semiample},\qquad R_1=N(J_1+bC_1)\geq0,
\end{equation}
with rational terms. On this fixed model choose a fresh general
fractional representative of \(bC_1\), so that its sum with \(B_1\)
is again simple normal crossing and klt.
Apply Proposition~\ref{pre:contract} to this augmented ordinary
adjoint and its known negative part in
\eqref{finish:augmented-decomposition}. Its finite ordinary scaling
contracts or flips detected analytic extremal rays negative for the
augmented adjoint. It preserves the actual semiample positive line
\(P_1\), with a fixed generated Cartier multiple, and contracts the
negative part \(R_1\).

The positive line \(P_1\) and the fixed algebraic-base line
\(C_1=h_1^*A_W\) have distinct roles; they need not be equal or
proportional. Preservation of \(C_1\) and of the specific morphism
to \(W\) follows from a separate ray estimate. We check inductively
that each step is over \(W\). Suppose the current
model \(T_i\) still has a morphism \(h_i:T_i\to W\) and the unaugmented
pair is klt. Put
\[
 J_i=K_{T_i}+B_i,\qquad C_i=h_i^*A_W.
\]
An extremal ray negative for \(J_i+bC_i\) is \(J_i\)-negative because
\(C_i\) is nef. The K\"ahler klt cone theorem
\cite[Theorem~1.3]{HX}, with zero nef b-part, supplies a rational curve
\(\Gamma\) generating this ray and satisfying
\[
 0<-J_i\cdot\Gamma\leq2n.
\]
If \(C_i\cdot\Gamma>0\), Cartier integrality gives
\[
 (J_i+bC_i)\cdot\Gamma\geq-2n+b>0,
\]
a contradiction. Hence \(C_i\) is numerically trivial on the contracted
ray.

The contracted fibers are projective, and all their curves have class
in the contracted ray. If their image under \(h_i\) had positive
dimension, a curve on a fiber would have positive degree against
\(h_i^*A_W\); thus \(h_i\) is constant on each contracted fiber.
Connected-fiber descent to the normal contraction base factors \(h_i\)
through it. For a flip, composing the morphism from the flipped space
with this base map gives the next \(h_i\). The added line is consequently
the same Cartier pullback \(C_{i+1}=h_{i+1}^*A_W\) on the next model.
Each step is also \(J_i\)-negative, and the unaugmented ordinary klt
pair persists. This proves the induction through the finite program.
On a common resolution the two pullbacks of the line from \(W\)
agree. Subtracting \(b\) times this pullback from the augmented
comparison therefore gives an effective exceptional comparison for
the unaugmented adjoints as well.
Writing \(P_Y\) for the descended positive line, the final actual
rational-line identity is
\[
 K_Y+\Delta+b h_Y^*A_W\qsim P_Y.
\]
The line \(P_Y\) is semiample by Proposition~\ref{pre:contract}, which
proves the required semiampleness of the augmented adjoint.
\end{proof}

The degree argument is the familiar way of preserving a nef Cartier
line in an adjoint program; compare \cite[Lemma~3.6]{K} and
\cite[Section~8.2]{P}. Here preservation of the line also preserves the
specific morphism to \(W\).

\subsection{The strengthened decomposition}

\begin{proof}[Proof of Theorem~\ref{thm:decomposition}]
We induct simultaneously on \(n=\dim T\) for both assertions of the
theorem. Dimension zero is immediate. Suppose both assertions hold in
every dimension less than \(n\). The projective case of the decomposition
is Proposition~\ref{pre:projective}. If \(a(T)=0\),
Proposition~\ref{zero:vanishing} gives \(M\num0\); ordinary decomposition
then proves the required statement for \(J+M\). It remains to treat
\[
 0<a(T)<n.
\]

Resolve the algebraic reduction. This changes the ordinary adjoint by
\eqref{pre:resolution-error}, so we may work on the resulting smooth
K\"ahler model and remove its exceptional error at the end.
We obtain a morphism \(h:T\to W\) with connected fibers and smooth
projective \(W\), where \(\dim W=a(T)\).
By Proposition~\ref{cycle:nef-summand} and
Corollary~\ref{cycle:base-descent}, after further modifications if
needed, there is a nef rational line \(N_W\) on \(W\) such that
\begin{equation}\label{finish:nef-descent}
 M\num h^*N_W.
\end{equation}
Apply Proposition~\ref{finish:over-base}. On its final model \(Y\),
write \(J_Y=K_Y+\Delta\). The map defined by a generated multiple of
\(J_Y+b h_Y^*A_W\), together with \(h_Y\), has image in a product of
projective varieties. Its Stein factorization is
\[
 f:Y\longrightarrow Z,\qquad j:Z\longrightarrow W,
\]
where \(Z\) is normal projective, \(f\) has connected fibers, and
\(h_Y=j\circ f\). The semiample augmented line comes from \(Z\);
subtracting \(b j^*A_W\) gives a rational Cartier divisor \(H\) with
\begin{equation}\label{finish:adjoint-base}
 J_Y\qsim f^*H.
\end{equation}
Surjectivity to \(W\) and the definition of algebraic dimension give
\[
 a(T)=\dim W\leq\dim Z\leq a(Y)=a(T).
\]
In particular \(f\) has positive relative dimension.

A divisible multiple of \(J_1\) has a nonzero section by
Theorem~\ref{pre:ordinary}. The effective exceptional comparisons in
the ordinary negative program preserve these section spaces.
Equation~\eqref{finish:adjoint-base} and \(f_*\cO_Y=\cO_Z\) therefore
give a nonzero section of a multiple of \(H\); in particular \(H\)
is pseudo-effective. Proposition~\ref{adj:base} gives
\[
 H\qsim K_Z+\Delta_Z
\]
for an effective rational boundary \(\Delta_Z\) with \((Z,\Delta_Z)\)
klt. The divisor
\[
 H+j^*N_W
\]
now satisfies exactly the hypotheses of Proposition~\ref{pre:projective}.

Accordingly, take a smooth projective modification \(u:Z'\to Z\)
and a birational morphism \(v:Z'\to Z_m\) such that
\begin{equation}\label{finish:projective-decomposition}
 u^*(H+j^*N_W)\qsim v^*H_m+E,
\end{equation}
where \(H_m\) is nef rational Cartier, \(E\geq0\) is rational and
\(v\)-exceptional, and \(v^*H_m\) is numerically equivalent to a
semiample rational line \(P_{Z'}\).

Take a smooth compact K\"ahler common resolution \(V\) of the program
and of the main transform over \(Z'\). Denote its morphisms by
\[
 q:V\to T_1,\qquad p:V\to Y,\qquad g:V\to Z',
 \qquad r=v\circ g.
\]
Thus \(f\circ p=u\circ g\). Pulling back
\eqref{finish:projective-decomposition} gives
\[
 p^*(J_Y+h_Y^*N_W)\qsim r^*H_m+g^*E.
\]
Every component of \(g^*E\) maps into a subset of codimension at least
two in \(Z_m\). Lemma~\ref{finish:exceptional} therefore identifies
\[
 N\bigl(p^*(J_Y+h_Y^*N_W)\bigr)=g^*E.
\]
Its positive class is represented by the semiample rational line
\(g^*P_{Z'}\).

Finally, the ordinary program comparison gives
\[
 q^*J_1\qsim p^*J_Y+F,\qquad F\geq0
 \text{ exceptional over }Y.
\]
The line from \(W\) is unchanged throughout the program. Combining
this equality with \eqref{finish:nef-descent} yields
\[
 q^*(J_1+M_1)\num g^*P_{Z'}+g^*E+F,
\]
where \(M_1\) is the pullback of \(M\). By exceptional translation,
\[
 N\bigl(q^*(J_1+M_1)\bigr)=g^*E+F.
\]
This is the required decomposition on the model \(T_1\).
Removing the effective resolution errors by
Lemma~\ref{pre:negative}(5) proves it for the original \(T\).
The zero-algebraic-dimension assertion and the decomposition have now
both been established in dimension \(n\), completing the simultaneous
induction.
\end{proof}

\subsection{Descent of the semiample representative}

\begin{proof}[Proof of Theorem~\ref{thm:main}]
Apply Theorem~\ref{thm:decomposition} to
\(D=K_X+B+M\). It gives a smooth compact K\"ahler modification
\(\mu:U\to X\) and
\[
 \mu^*D\num P+R,\qquad P\text{ semiample},\qquad
 R=N(\mu^*D).
\]
Since \(D\) is nef, \(\mu^*D\) is nef and \(R=0\).
Thus \(P-\mu^*D\) has zero real Chern class. By
Lemma~\ref{pre:picard}, it is the pullback of a rational line
\(F\in\Pic^0(X)\otimes_{\Z}\Q\).
Set \(L=D+F\). Then \(c_1(L)=c_1(D)\) and
\(\mu^*L\qsim P\). The last assertion of Lemma~\ref{pre:picard}
descends semiampleness to \(L\).
\end{proof}
